On prépare une solution aqueuse en dissolvant une masse $m=\qty{277.5}{\mg}$ de
chlorure de calcium $\ce{CaCl2}$ dans un volume $V=\qty{250}{\mL}$ d'eau
distillée.

\begin{donnees}
    \begin{itemize}
        \item $M(\ce{CaCl2})=\qty{111.1}{\g\per\mol}$
        \item Les conductivités molaires ioniques~: \\
              $\lambda(\ce{Cl^-})=\qty{7.63e-3}{\lambda}$ et
              $\lambda(\ce{Ca^2+})=\qty{1.19e-2}{\lambda}$.
    \end{itemize}
\end{donnees}

\begin{enumerate}
    \item Écrire l'équation de dissolution de $\ce{CaCl2}$ dans
          l'eau.~\textbf{(0.5pt)}
    \item Déterminer la concentration $C$ de la solution obtenue en
          $\unit{\mol\per\L}$ puis en $\unit{\mol\per\cubic\m}$.~\textbf{(1pt)}

    \item Déterminer les concentrations effectives des espèces chimiques qui se
          trouvent dans la solution.~\textbf{(1.5pt)}
    \item Déterminer la conductivité $\sigma$ de la solution.~\textbf{(1pt)}
\end{enumerate}
